%% appendix/datasheet.tex — Gebru et al. (2021) Datasheet for Datasets
%% Filled in for the MIRROR benchmark and dataset release.
%% Template follows: Gebru et al., "Datasheets for Datasets," CACM 2021.

\subsection*{Overview}

This datasheet covers the \mirror{} benchmark in its entirety, including: (1) the
5,000-question base question bank (Experiments 1--5), (2) the Experiment 9 composite
task bank (597 tasks), (3) the trial result files from all experiments, and (4) the
analysis code and scripts. We follow the Gebru \etal{} datasheet template~\citep{gebru2021datasheets}.

% ─────────────────────────────────────────────────────────────────────────────
\subsection*{Motivation}
% ─────────────────────────────────────────────────────────────────────────────

\textbf{For what purpose was the dataset created?}
\mirror{} was created to provide the first systematic, hierarchical benchmark for
metacognitive calibration in large language models. The benchmark operationalizes
four ascending levels of metacognitive capability---self-knowledge, knowledge transfer,
compositional self-prediction, and adaptive self-regulation---and connects them to
behavioral outcomes in agentic task settings. It was created to support the NeurIPS
2026 paper ``\mirror{}: A Hierarchical Benchmark for Metacognitive Calibration in
Large Language Models'' and subsequent research.

\textbf{Who created the dataset and on behalf of which entity?}
The dataset was created by the \mirror{} project team. Details provided after review.

\textbf{Who funded the creation of the dataset?}
Details provided after review. API costs were funded by research compute credits.

% ─────────────────────────────────────────────────────────────────────────────
\subsection*{Composition}
% ─────────────────────────────────────────────────────────────────────────────

\textbf{What do the instances that comprise the dataset represent?}
The dataset consists of three artifact types:

\begin{enumerate}
  \item \textbf{Question bank} (\texttt{data/questions/}): 5,000 multiple-choice questions
    (625 per domain $\times$ 8 domains). Each question has a unique identifier, domain label,
    subcategory label (5 per domain, 40 total), difficulty rating, question text, 4 answer
    choices, and a verified correct answer. Questions were generated by frontier LLMs and
    verified by cross-model consensus.

  \item \textbf{Experiment 9 task bank} (\texttt{data/exp9\_tasks.jsonl}): 597 composite
    tasks (297 fixed, 300 tailored). Each task has the following schema:
    \begin{itemize}
      \item \texttt{task\_id} (string): unique identifier
      \item \texttt{task\_type} (``fixed'' | ``tailored''): fixed = identical across all models;
        tailored = model-specific strong/weak domain pair
      \item \texttt{circularity\_free} (bool): True for fixed tasks (primary analysis)
      \item \texttt{target\_model} (string | null): model slug for tailored tasks
      \item \texttt{domain\_a}, \texttt{domain\_b} (string): one of 8 domains
      \item \texttt{subcategory\_a}, \texttt{subcategory\_b} (string): one of 40 subcategories
      \item \texttt{difficulty\_a}, \texttt{difficulty\_b} (``easy''|``medium''|``hard'')
      \item \texttt{correct\_answer\_a}, \texttt{correct\_answer\_b} (string)
      \item \texttt{answer\_type\_a}, \texttt{answer\_type\_b} (``multiple\_choice''|``short\_answer'')
      \item \texttt{task\_text} (string): combined task description
      \item \texttt{part1\_text}, \texttt{part2\_text} (string): per-component prompts
    \end{itemize}

  \item \textbf{Trial result files} (\texttt{data/results/}): JSONL files with one record
    per model per task per condition per paradigm. Trial record schema (Experiment 9):
    \begin{itemize}
      \item \texttt{model} (string), \texttt{task\_id} (string)
      \item \texttt{condition} (int, 1--4), \texttt{paradigm} (int, 1--3)
      \item \texttt{is\_false\_score\_control} (bool), \texttt{circularity\_free} (bool)
      \item \texttt{domain\_a}, \texttt{domain\_b}, \texttt{subcategory\_a}, \texttt{subcategory\_b}
      \item \texttt{strength\_a}, \texttt{strength\_b} (``strong''|``weak''|``medium''|``unknown'')
      \item \texttt{component\_a\_decision}, \texttt{component\_b\_decision}
        (``proceed''|``use\_tool''|``defer'')
      \item \texttt{component\_a\_correct}, \texttt{component\_b\_correct} (bool)
      \item \texttt{component\_a\_answer}, \texttt{component\_b\_answer} (string)
      \item \texttt{exp1\_accuracy\_a}, \texttt{exp1\_accuracy\_b} (float, 0--1)
      \item \texttt{mirror\_gap\_a}, \texttt{mirror\_gap\_b} (float)
      \item \texttt{hedge\_count}, \texttt{decomp\_count}, \texttt{token\_count},
        \texttt{error\_type} (Paradigm 3 only)
    \end{itemize}
\end{enumerate}

\textbf{How many instances are there in total?}
\begin{itemize}
  \item Question bank: 5,000 questions
  \item Experiment 9 task bank: 597 tasks
  \item Experiment 9 trial records: 39,150 valid trials (from 46,914 total, before
    exclusion of failed API calls and excluded models)
  \item Across all experiments (Exp.\ 1--6, 8--9): $>$100,000 total trial records
\end{itemize}

\textbf{Does the dataset contain all possible instances or is it a sample?}
The question bank is a sample from the conceptual space of all valid questions
in each domain. The Experiment 9 task bank is a sample from the space of valid
composite tasks (two-domain pairs with verified correct answers). The trial records
are a complete enumeration of the experimental runs performed.

\textbf{What data does each instance consist of?}
See schema descriptions above. All data is text-based (question text, answer text,
model responses). No images, audio, or video. Model response text is the raw
output from the API calls and may contain model-generated content of any type.

\textbf{Is there a label or target associated with each instance?}
Yes. Each question has a verified correct answer. Each trial record has ground-truth
correct answer fields and model decision fields (computed by the response parser).

\textbf{Is any information missing from individual instances?}
Some trial records have null values for fields that could not be parsed from the
model's response (e.g., \texttt{component\_a\_answer} = ``'' if parsing failed).
Paradigm 3 behavioral fields (hedge\_count, etc.) are null for Paradigm 1 and 2
records. Approximately 3--4\% of trial records for non-excluded models have
\texttt{api\_failure} flags.

\textbf{Are there any relationships between individual instances?}
Yes. Questions within the question bank are grouped by domain and subcategory.
Experiment 9 tasks pair questions from two different domains. Trial records share
task\_id and model identifiers across conditions and paradigms.

\textbf{Are there any recommended data splits?}
\begin{itemize}
  \item Primary analysis (Experiment 9): use \texttt{circularity\_free=True} tasks
    (297 fixed tasks) as the primary split; \texttt{circularity\_free=False} tasks
    (300 tailored) as supplementary.
  \item Experiment 1: no train/test split; the question bank is used for evaluation only.
  \item Cross-model comparison: use models with $>0\%$ API success rate (i.e., exclude
    qwen-3-235b and command-r-plus from Experiment 9 analyses).
\end{itemize}

\textbf{Are there any errors, sources of noise, or redundancies in the dataset?}
\begin{itemize}
  \item Question bank: questions generated by LLMs may contain errors. Cross-model
    consensus verification reduces but does not eliminate incorrect ``correct answers.''
    Users should independently verify a sample before using the bank for high-stakes
    evaluation.
  \item Task bank: all 597 tasks pass static schema verification. LLM-based
    correctness verification (cross-model consensus check) was not completed before
    submission; some tasks may have incorrect correct\_answer fields.
  \item Trial records: model responses are raw API output. The response parser may
    misclassify decisions (proceed/use\_tool/defer) in edge cases where the model
    does not follow the prompt format. Parse failures are flagged with
    \texttt{parse\_success=False}.
\end{itemize}

\textbf{Is the dataset self-contained, or does it link to or otherwise rely on
external resources?}
Self-contained. All text data is included. The dataset does not include model weights
or API keys. API access to the evaluated models is required to reproduce the
experimental runs; the result files are included so that reproduction of the analysis
does not require re-running the API calls.

\textbf{Does the dataset contain data that might be considered confidential?}
No. The dataset contains only question-answer pairs in public knowledge domains and
model API responses. No personal data is included.

\textbf{Does the dataset contain data that, if viewed directly, might be offensive,
insulting, threatening, or might otherwise cause anxiety?}
The question bank domains (arithmetic, spatial, temporal, linguistic, logical, social,
factual, procedural) do not target sensitive topics. Social domain questions involve
theory-of-mind and norm-violation scenarios that are based on published benchmarks
(Social IQa, FOLIO). A small number of questions may describe hypothetical conflict
or norm violation scenarios in abstract terms. These are not expected to cause
distress to dataset users.

% ─────────────────────────────────────────────────────────────────────────────
\subsection*{Collection Process}
% ─────────────────────────────────────────────────────────────────────────────

\textbf{How was the data associated with each instance acquired?}
\begin{itemize}
  \item Question bank: generated by prompting frontier LLMs (GPT-4o, Claude-3.5-sonnet,
    Gemini-1.5-pro) with domain- and subcategory-specific generation prompts.
    Verified by cross-model consensus: a question is retained only if at least 2 of 3
    verifier models identify the same correct answer. The generation and verification
    scripts are released in \texttt{scripts/}.
  \item Experiment 9 task bank: generated by a template library
    (\texttt{mirror/data/exp9\_template\_library.py}, 37 pairs $\times$ 5 = 185 templates)
    combined with LLM-based task instantiation.
  \item Trial records: collected via API calls to model providers (NVIDIA NIM, DeepSeek,
    Google Vertex, Groq) using the unified async client (\texttt{mirror/api/client.py}).
\end{itemize}

\textbf{What mechanisms or procedures were used to collect the data?}
All data collection uses asynchronous Python scripts with exponential backoff retry,
JSONL output with fsync for crash resistance, and resume support. Concurrency is
32 parallel API calls per model for Experiment 9. All scripts are released in
\texttt{scripts/}.

\textbf{If the dataset is a sample of a larger set, what was the sampling strategy?}
The question bank is sampled by domain and subcategory to ensure even coverage.
Experiment 9 fixed tasks are sampled from the template library to cover all 8 domains
and 40 subcategories. Tailored tasks are sampled to match each model's specific
strong/weak domain profile from Experiment 1.

\textbf{Who was involved in the data collection process and how were they compensated?}
Data collection was performed by the research team using automated API calls. No crowd
workers were used. Model API costs were covered by research compute credits.

\textbf{Over what timeframe was the data collected?}
Experiments were conducted between January and March 2026. Experiment 9 full run:
\texttt{run\_id=20260312T140842} (run started March 12, 2026).

\textbf{Were any ethical review processes conducted?}
No IRB review was required; the dataset does not involve human subjects. The
benchmark evaluates AI systems, not human participants.

% ─────────────────────────────────────────────────────────────────────────────
\subsection*{Preprocessing / Cleaning / Labeling}
% ─────────────────────────────────────────────────────────────────────────────

\textbf{Was any preprocessing/cleaning/labeling of the data done?}
Yes. Question bank preprocessing includes:
\begin{itemize}
  \item Deduplication by question text hash (removing near-duplicates with $>$95\%
    string similarity; see \texttt{mirror/data/deduplicator.py}).
  \item Difficulty validation: questions are classified as easy/medium/hard by
    a difficulty validator model (\texttt{mirror/data/difficulty\_validator.py}).
  \item Answer normalization: correct answers are normalized to a canonical form
    (e.g., ``A'', ``B'', ``C'', ``D'' for multiple-choice) by the answer matcher
    (\texttt{mirror/scoring/answer\_matcher.py}).
\end{itemize}

Trial records preprocessing includes:
\begin{itemize}
  \item Response parsing: raw API response text is parsed to extract decisions,
    answers, and behavioral signals by paradigm-specific parsers.
  \item Answer matching: model answers are matched to correct answers using a robust
    fuzzy matcher (\texttt{mirror/scoring/answer\_matcher.py},
    \texttt{match\_answer\_robust}).
  \item Exclusion tagging: records from excluded models (qwen-3-235b, command-r-plus)
    are tagged but retained in the raw JSONL for completeness.
\end{itemize}

\textbf{Was the ``raw'' data saved in addition to the preprocessed/cleaned/labeled data?}
Yes. Raw API response text is preserved in the \texttt{raw\_response} field of each
trial record. Raw question generation outputs are preserved in
\texttt{data/questions/raw/}.

\textbf{Is the software used to preprocess/clean/label the instances available?}
Yes. All preprocessing scripts are released in \texttt{scripts/} and
\texttt{mirror/data/}. Key files:
\begin{itemize}
  \item \texttt{mirror/data/deduplicator.py} — deduplication
  \item \texttt{mirror/data/difficulty\_validator.py} — difficulty classification
  \item \texttt{mirror/scoring/answer\_matcher.py} — answer matching
  \item \texttt{mirror/experiments/agentic\_paradigms.py} — response parsing
\end{itemize}

% ─────────────────────────────────────────────────────────────────────────────
\subsection*{Uses}
% ─────────────────────────────────────────────────────────────────────────────

\textbf{Has the dataset been used for any tasks already?}
Yes. The dataset has been used to produce all results reported in the accompanying
NeurIPS 2026 paper. Pre-registered analysis plans are archived at \todo{OSF URL}.

\textbf{Is there a repository that links to any or all papers or systems that use the dataset?}
The repository is at \todo{GitHub URL}. All paper code and data are released under
the MIT License (code) and CC BY 4.0 (data).

\textbf{What (other) tasks could the dataset be used for?}
\begin{itemize}
  \item Calibration benchmarking: the 5,000-question bank with domain labels can serve
    as a calibration evaluation set for any new LLM.
  \item Agentic evaluation: the Experiment 9 task bank can evaluate tool-use and
    deferral decisions in new agentic systems without repeating the full \mirror{} protocol.
  \item Metacognitive research: the trial records provide a rich dataset for studying
    the relationship between confidence expression and task performance in LLMs.
  \item Sycophancy analysis: Experiment 6 data can be used to study social pressure
    sensitivity across model families.
\end{itemize}

\textbf{Is there anything about the composition of the dataset or the way it was
collected and preprocessed/cleaned/labeled that might impact future uses?}
\begin{itemize}
  \item The question bank was generated by LLMs; it will not capture genuinely novel
    question types that are hard for current LLMs to generate correctly.
  \item The wagering channel uses a binary (high/low) format; this is not equivalent
    to continuous probability elicitation.
  \item Model API results are specific to the model versions tested in 2026. Results
    may not generalize to updated model versions.
  \item The circularity problem: Experiment 9 fixed tasks (primary analysis) were
    generated without knowledge of specific model accuracy profiles, avoiding the
    circularity of tailored tasks. Users should use fixed tasks for primary comparisons
    and tailored tasks only as supplementary.
\end{itemize}

\textbf{Are there tasks for which the dataset should not be used?}
The dataset should not be used as a general-purpose question-answering benchmark
(the questions are designed for calibration measurement, not capability evaluation
in isolation). The trial records include model API responses and should not be used
to train models that would then be evaluated on the same question bank (data
contamination concern).

% ─────────────────────────────────────────────────────────────────────────────
\subsection*{Distribution}
% ─────────────────────────────────────────────────────────────────────────────

\textbf{Will the dataset be distributed to third parties outside of the entity on
behalf of which the dataset was created?}
Yes. The full dataset will be publicly released upon paper acceptance.

\textbf{How will the dataset be distributed?}
Via GitHub (\todo{URL}) and Hugging Face Datasets (\todo{HF URL}).
The question bank and task bank are distributed as JSONL files.
Trial records are distributed as compressed JSONL archives (one per experiment).

\textbf{When will the dataset be distributed?}
Upon NeurIPS 2026 acceptance or camera-ready deadline, whichever comes first.

\textbf{Will the dataset be distributed under a copyright or other intellectual
property (IP) license, and if so, which one?}
Data: Creative Commons Attribution 4.0 International (CC BY 4.0).
Code: MIT License.

\textbf{Have any third parties imposed IP-based or other restrictions on the data
associated with the instances?}
No. The questions are generated by LLMs; no third-party IP restrictions apply.
Some questions are inspired by existing benchmarks (ARC, GSM8k, TriviaQA, Social IQa)
but are newly generated instances, not copies of those benchmark items.

\textbf{Do any export controls or other regulatory restrictions apply to the dataset or
to individual instances?}
No.

% ─────────────────────────────────────────────────────────────────────────────
\subsection*{Maintenance}
% ─────────────────────────────────────────────────────────────────────────────

\textbf{Who will be supporting/hosting/maintaining the dataset?}
The research team. Contact details provided after review.

\textbf{How can the owner/curator/manager of the dataset be contacted?}
Via the GitHub repository issue tracker (\todo{URL}).

\textbf{Is there an erratum?}
Not at time of initial release. An erratum section will be maintained in the
GitHub repository README.

\textbf{Will the dataset be updated?}
Yes. Planned updates:
\begin{itemize}
  \item Addition of 270 tailored tasks for the supplementary analysis (blocked on
    Experiment 1 backfill completion for deepseek-v3, phi-4, command-r-plus).
  \item Experiment 5 ARS completion for gemini-2.5-pro and qwen-3-235b
    (analysis scripts ready; raw data collected).
  \item Paradigm 3 behavioral signal analysis for Experiment 9.
\end{itemize}
Updates will be versioned (v1.0 = initial release, v1.1 = first update, etc.).

\textbf{If the dataset relates to people, are there applicable limits on the
retention of the data associated with the instances?}
Not applicable. The dataset does not contain personal data.

\textbf{Will older versions of the dataset continue to be supported/hosted/maintained?}
Yes. All versions will be tagged in the GitHub repository and archived on Zenodo.

\textbf{If others want to extend/augment/build on/contribute to the dataset, is
there a mechanism for them to do so?}
Yes. Contributions via GitHub pull requests are welcome. Guidelines for contributing
new questions (format, verification requirements) will be provided in the repository
CONTRIBUTING.md.
